فرض کنید $p_{k}$، $k$-امین عدد اول باشد برای $k \geq 1$. قرار دهید $a_{1} = 2$. برای $n \geq 1$، $a_{n+1}$ را کوچک‌ترین عدد صحیح بزرگ‌تر از $a_{n}$ بگیرید که به ازای تمام $k \leq n$، به پیمانه‌ی $p_{k+1}$ هم‌نهشت با $-k$ باشد. چنین عددی بنابر قضیه‌ی باقیمانده‌ی چینی وجود دارد.

بنابراین، برای هر $k \geq 0$، داریم $k + a_{n} \equiv 0 \pmod{p_{k+1}}$ هرگاه $n \geq k+1$. در نتیجه، حداکثر $k+1$ جمله از دنباله‌ی $\{k + a_{n}\}$ می‌تواند اول باشد: از جمله‌ی $(k+2)$-ام به بعد، هر مقدار مضرب غیربدیهی از $p_{k+1}$ است و در نتیجه مرکب است.

این اثبات را کامل می‌کند.
